\documentclass[a4paper,11pt]{article}
\usepackage[T1]{fontenc}
\usepackage[utf8]{inputenc}
\usepackage[italian]{babel}
\usepackage{amsmath,amssymb}
\usepackage{graphicx}
\usepackage{geometry}
\usepackage{xcolor}
\usepackage{listings}
\geometry{margin=2.5cm}

\lstdefinelanguage{MiniZinc}{
  morekeywords={int,set,of,array,var,constraint,include,solve,satisfy,output,show},
  sensitive=true,
  morecomment=[l]{\%},
  morestring=[b]"
}

\lstset{
  language=MiniZinc,
  basicstyle=\ttfamily\small,
  keywordstyle=\color{blue}\bfseries,
  commentstyle=\color{teal!70!black}\itshape,
  stringstyle=\color{orange!70!black},
  numbers=left,
  numberstyle=\tiny\color{gray},
  stepnumber=1,
  numbersep=8pt,
  showstringspaces=false,
  frame=single,
  breaklines=true
}

\title{FOL: Aldo}
\date{}

\begin{document}
\maketitle

\section*{Problema}
Modellare mediante una formula $\phi$ in logica del primo ordine le seguenti affermazioni:
\begin{quote}
\textit{I cani da caccia abbaiano di notte. Chi possiede gatti non può avere topi in casa.
Chi ha il sonno leggero non può possedere nessun animale che abbaia di notte.
Aldo possiede un gatto oppure un cane da caccia.}
\end{quote}

\subsection*{1.1 Modellazione}

\paragraph{Costanti e Predicati}
\begin{itemize}
  \item $Aldo$: Aldo
  \item $Cane(x)$: $x$ è un cane da caccia
  \item $Gatto(x)$: $x$ è un gatto
  \item $Topo(x)$: $x$ è un topo
  \item $Animale(x)$: $x$ è un animale
  \item $Abbaia(x)$: $x$ abbaia di notte
  \item $SonnoLeggero(x)$: $x$ ha il sonno leggero
  \item $Possiede(x, y)$: $x$ possiede (o ha in casa) $y$
\end{itemize}

\paragraph{Formula}
La formula $\phi$ è data dalla congiunzione delle singole affermazioni:
\[
\phi \equiv \phi_1 \land \phi_2 \land \phi_3 \land \phi_4 \land \phi_5
\]
dove:
    
I cani da caccia abbaiano di notte:
\[\phi_1 = \forall x \, (Cane(x) \Rightarrow Abbaia(x))\]

Chi possiede gatti non può avere topi in casa:
\[\phi_2 = \forall x \, ((\exists y \, (Gatto(y) \land Possiede(x, y))) \Rightarrow \neg \exists z \, (Topo(z) \land Possiede(x, z)))\]

Chi ha il sonno leggero non può possedere nessun animale che abbaia di notte:
\[\phi_3 = \forall x \, (SonnoLeggero(x) \Rightarrow \neg \exists y \, (Animale(y) \land Abbaia(y) \land Possiede(x, y)))\]

Aldo possiede un gatto oppure un cane da caccia:
\[\phi_4 = (\exists x \, (Gatto(x) \land Possiede(Aldo, x))) \lor (\exists y \, (Cane(y) \land Possiede(Aldo, y)))\]

Tutti i cani da caccia e i gatti sono animali:
\[\phi_5 = \forall x \, (Cane(x) \Rightarrow Animale(x)) \land \forall x \, (Gatto(x) \Rightarrow Animale(x))\]


\newpage

\subsection*{1.2 Modello $M$ di $\phi$}

Costruiamo un interpretazione $M$ che renda vera la formula $\phi$.

\paragraph{Dominio e interpretazione}
Scegliamo il dominio $D^M = \{\mathrm{Aldo}, \mathrm{Spillo}\}$.
\begin{itemize}
  \item $Aldo^M = \mathrm{Aldo}$
  \item $Cane^M = \varnothing$
  \item $Gatto^M = \{\mathrm{Spillo}\}$
  \item $Topo^M = \varnothing$
  \item $Animale^M = \{\mathrm{Spillo}\}$
  \item $Abbaia^M = \varnothing$
  \item $SonnoLeggero^M = \varnothing$
  \item $Possiede^M = \{(\mathrm{Aldo}, \mathrm{Spillo})\}$
\end{itemize}

\paragraph{Verifica che $M \models \phi$}
\begin{itemize}
    \item $\phi_1$ è vera perché $Cane^M = \varnothing$ (antecedente falso).
    \item $\phi_2$ è vera: Aldo possiede un gatto, ma poiché non esistono topi nel dominio ($Topo^M = \varnothing$), la negazione dell esistenza di topi posseduti è vera.
    \item $\phi_3$ è vera perché $SonnoLeggero^M = \varnothing$ (antecedente falso).
    \item $\phi_4$ è vera: Aldo possiede $\mathrm{Spillo}$, che è un gatto. Dunque la prima parte dell'OR è soddisfatta.
    \item $\phi_5$ è vera: l'unico gatto appartiene ad $Animale^M$.
\end{itemize}
Poiché tutte le sottoformule sono vere, l'intera congiunzione $\phi$ è vera in $M$. Quindi, $M \models \phi$.

\newpage

\subsection*{1.3 Non-modello $I$ di $\phi$}

Costruiamo un interpretazione $I$ in cui $\phi$ sia falsa, falsificando $\phi_2$.

\paragraph{Dominio e interpretazione}
Scegliamo il dominio $D^I = \{\mathrm{Aldo}, \mathrm{Spillo}, \mathrm{Topolino}\}$.
\begin{itemize}
  \item $Aldo^I = \mathrm{Aldo}$
  \item $Cane^I = \varnothing$
  \item $Gatto^I = \{\mathrm{Spillo}\}$
  \item $Topo^I = \{\mathrm{Topolino}\}$
  \item $Animale^I = \{\mathrm{Spillo}\}$
  \item $Abbaia^I = \varnothing$
  \item $SonnoLeggero^I = \varnothing$
  \item $Possiede^I = \{(\mathrm{Aldo}, \mathrm{Spillo}), (\mathrm{Aldo}, \mathrm{Topolino})\}$
\end{itemize}

\paragraph{Verifica che $I \not\models \phi$}
Analizziamo la formula $\phi_2$:
\[
\forall x \, ((\exists y \, (Gatto(y) \land Possiede(x, y))) \Rightarrow \neg \exists z \, (Topo(z) \land Possiede(x, z)))
\]
Per $x = \mathrm{Aldo}$:
\begin{itemize}
    \item L'antecedente è vero per $y = \mathrm{Spillo}$.
    \item Il conseguente richiede $\neg \exists z \, (Topo(z) \land Possiede(\mathrm{Aldo}, z))$.
    \item Tuttavia, per $z = \mathrm{Topolino}$, $Topo^I(\mathrm{Topolino})$ e $Possiede^I(\mathrm{Aldo}, \mathrm{Topolino})$ sono veri.
    \item Quindi l'esistenza del topo posseduto da Aldo rende il conseguente falso.
\end{itemize}
L'implicazione in $\phi_2$ risulta falsa (Vero $\Rightarrow$ Falso). Dunque $\phi$ complessiva è falsa e $I \not\models \phi$.

\end{document}